\documentclass[11pt]{article}
\usepackage[T1]{fontenc}
\usepackage{amsmath}
\usepackage{amssymb}
\usepackage{amsthm}
\usepackage[hidelinks]{hyperref}
\newtheorem{theorem}{Theorem}[section]
\newtheorem{lemma}[theorem]{Lemma}
\newtheorem{corollary}[theorem]{Corollary}
\theoremstyle{definition}
\newtheorem{definition}[theorem]{Definition}
\begin{document}
\begin{center}
{\LARGE Natural number divisibility}
\end{center}
\section{Divisibility}
\label{ll:divisibility:definition}
\begin{definition}
\label{ll:divisibility:divides}
For every natural number \(n\) and natural number \(m\), \(\texttt{peano.arith::divides}\) holds exactly when there exists a natural number \(k\) such that \(m = n \cdot k\).
\end{definition}
\section{Divisibility : order}
\label{ll:divisibility:properties}
\begin{theorem}
\label{ll:divisibility:divides-refl}
For every natural number \(n\), \(\operatorname{divides}(n, n)\).
\end{theorem}
\begin{proof}
Use \(1\) as the witness.
The goal follows from \(\operatorname{eqsymm}(\operatorname{mulone}(n))\).
\end{proof}
\begin{theorem}
\label{ll:divisibility:one-divides}
For every natural number \(n\), \(\operatorname{divides}(1, n)\).
\end{theorem}
\begin{proof}
Use \(n\) as the witness.
The goal follows from \(\operatorname{eqsymm}(\operatorname{onemul}(n))\).
\end{proof}
\begin{theorem}
\label{ll:divisibility:divides-trans}
For every natural number \(a\) and natural number \(b\) and natural number \(c\), if \(\operatorname{divides}(a, b)\) and \(\operatorname{divides}(b, c)\), then \(\operatorname{divides}(a, c)\).
\end{theorem}
\begin{proof}
Assume \(h\).
Consider the cases of \(h\).
Case \(andintro\) with \(hab\), \(hbc\):
Consider the cases of \(hab\).
Case \(existsintro\) with \(k\), \(hk\):
Consider the cases of \(hbc\).
Case \(existsintro\) with \(l\), \(hl\):
Use \(k \cdot l\) as the witness.
Rewrite the goal using \(\rightarrow hl\), \(\rightarrow hk\).
The goal follows from \(\operatorname{mulassoc}(a, k, l)\).
\end{proof}
\end{document}
